\uuid{esQK}
\titre{Composée/DL usuels}
\auteur{hanani ; mzari}
\organisation{exo7}
\licence{CC BY-NC-SA 4.0}
\datecreate{2020-01-20}
\niveau{L1}
\module{Analyse}
\chapitre{Développement limité}
\sousChapitre{Calculs}
\difficulte{1}
\theme{Opérations sur les DL, Composée, DL usuels}
\exo7qcmid{lille-2.developpements-limites.001}

\contenu{
\texte{
Soit $a\in \Rr$ et $f$ une fonction. On dit que $f(x)$ est un $o\left((x-a)^n\right)$ au voisinage de $a$ et on écrit $f(x)=o((x-a)^n)$ si $\displaystyle \lim _{x\to a}\frac{f(x)}{(x-a)^n}=0$. On note par $DL_n(a)f(x)$ le développement limité de $f$ à l'ordre $n$ en $a$. On note aussi $DL_n(a^+)f(x)$ (resp. $DL_n(a^-)f(x)$) le développement limité de $f$ à l'ordre $n$ à droite en $a$ (resp. à gauche en $a$).
}

\begin{question}
\qclassification{125.02}
\qtags{motcle=Composée/DL usuels}
Soit $f$ une fonction définie sur un intervalle ouvert contenant $0$ telle que $f(0)=0$ et $DL_4(0)f(x)=x+x^2+o(x^4)$. On peut en déduire que :
\begin{answers}
\good{$f$ est continue en $0$, dérivable en $0$ et $f'(0)=1$.}
\bad{Si $f$ est $2$ fois dérivable en $0$ alors $f^{(2)}(0)=1$.}
\bad{$\displaystyle DL_4(0)f(2x)=2x+2x^2+o(x^4)$.}
\good{$\displaystyle DL_4(0)f(x^2)=x^2+x^4+o(x^4)$.}
\end{answers}
\begin{explanations}
On a : $\displaystyle \lim _{x\to 0}f(x)=0=f(0)$. Donc $f$ est continue en $0$. L'existence du $DL$ à l'ordre $1$ en $0$ implique que $f$ est dérivable en $0$ et si $f$ est $2$ fois dérivable en $0$ alors $f^{(2)}(0)=2$. Enfin, $DL_2(0)f(2x)=2x+4x^2+o(x^4)$ et
$$DL_8(0)f(x^2)=x^2+x^4+o(x^8)\Rightarrow DL_4(0)f(x^2)=x^2+x^4+o(x^4).$$
\end{explanations}
\end{question}
}
